
% GA-PRAC-01
% C2 - Escenario 2: Sabores de la Abuela Rosa
% Responsable: Cordova Carreno Mayra Lucila
% Requiere: longtable y array

\subsection{Escenario 2: Sabores de la Abuela Rosa}

El emprendimiento Sabores de la Abuela Rosa necesita
una aplicación móvil para gestionar pedidos de comida
típica. El proyecto presenta alta volatilidad de
requisitos, disponibilidad frecuente de la propietaria,
un presupuesto pequeño y un plazo de 14 semanas.

El sistema debe proteger los datos personales de los
clientes, controlar los pagos por transferencia y
funcionar en teléfonos Android de gama baja con
conectividad intermitente.

\small
\setlength{\tabcolsep}{4pt}
\renewcommand{\arraystretch}{1.35}

\begin{longtable}{
|>{\raggedright\arraybackslash}p{0.16\textwidth}
|>{\raggedright\arraybackslash}p{0.25\textwidth}
|>{\raggedright\arraybackslash}p{0.25\textwidth}
|>{\raggedright\arraybackslash}p{0.25\textwidth}|}
\hline
\textbf{Criterio} &
\textbf{Secuencial (cascada)} &
\textbf{Iterativo e incremental} &
\textbf{Ágil} \\
\hline
\endfirsthead

\hline
\textbf{Criterio} &
\textbf{Secuencial (cascada)} &
\textbf{Iterativo e incremental} &
\textbf{Ágil} \\
\hline
\endhead

\textbf{K1. Estabilidad y cambios} &
Una especificación inicial quedaría desactualizada
por los cambios diarios del menú y las nuevas ideas
de Doña Rosa y su hijo. &
Permite incorporar progresivamente cambios en
promociones, menú y funcionalidades mediante
ciclos de revisión. &
Favorece la repriorización frecuente de requisitos
según los cambios del menú y las decisiones
de Doña Rosa. \\
\hline

\textbf{K2. Partes interesadas y disponibilidad} &
Doña Rosa podría aprobar requisitos al inicio,
pero su disponibilidad diaria quedaría poco
aprovechada durante el desarrollo. &
Permite que Doña Rosa pruebe los incrementos
y comunique correcciones al finalizar cada ciclo. &
Aprovecha su disponibilidad diaria para validar
funcionalidades y decidir prioridades, aunque
deben controlarse las solicitudes adicionales
de su hijo. \\
\hline

\textbf{K3. Tipo de sistema y criticidad} &
Permite documentar requisitos de pagos y protección
de datos, pero podría detectar tarde problemas
de uso en teléfonos económicos. &
Facilita probar progresivamente pedidos, pagos
y funcionamiento con conexiones intermitentes. &
Permite validar tempranamente la experiencia móvil,
pero exige criterios de aceptación para pagos,
privacidad y funcionamiento sin conexión estable. \\
\hline

\textbf{K4. Riesgo principal} &
Los cambios frecuentes podrían provocar retrabajo
y comprometer la entrega dentro de las 14 semanas. &
Existe riesgo de no completar los incrementos
previstos con dos desarrolladores a medio tiempo. &
Las nuevas ideas pueden ampliar excesivamente
el alcance y agotar el presupuesto antes
de las fiestas. \\
\hline

\textbf{K5. Documentación y marco normativo o contractual} &
Favorece una especificación formal de datos
personales y pagos, aunque requiere actualizarla
ante cambios frecuentes. &
Permite mantener requisitos y evidencias
de pruebas actualizados por incremento. &
Puede utilizar historias de usuario y criterios
de aceptación, complementados con documentación
de privacidad y conciliación de transferencias. \\
\hline

\end{longtable}

\normalsize
